\begin{recipebox}{Spicy Beef Bolognese with Rigatoni}{3--4}{40 min}
  \begin{ingredients}
    \item 500\,g beef mince
    \item 350--400\,g rigatoni, cooked separately
    \item 1 onion, diced
    \item 3 garlic cloves, minced
    \item 400\,g tinned tomatoes
    \item 2 tbsp tomato paste
    \item 1 tsp chilli flakes (adjust to taste)
    \item 1 tsp dried oregano (or Italian herbs)
    \item Splash red wine (optional)
    \item \nicefrac{1}{2}--1 cup beef stock or water, to loosen the sauce
    \item 250\,g frozen mixed vegetables
    \item Handful frozen spinach
    \item 1 tbsp olive oil
    \item Salt and black pepper, to taste
    \item Mozzarella, grated, to finish
    \item Parmesan, grated, to finish
  \end{ingredients}
  \begin{method}
    \item Heat the olive oil in a large pan or pot over medium-high heat. Add the mince and cook, breaking it up, until well browned all over (about 5 minutes).
    \item Add the onion and cook for 3--4 minutes until softened.
    \item Add the garlic, tomato paste, chilli flakes, and oregano. Stir and cook for 1 minute until fragrant.
    \item If using red wine, pour it in now and let it bubble and reduce for a couple of minutes, scraping up any browned bits.
    \item Add the tinned tomatoes along with the beef stock or water --- since the pasta is cooked separately this time rather than absorbing liquid in the pot, the sauce needs a bit more liquid than a standard bolognese to stay saucy. Season with salt and pepper.
    \item Bring to a simmer, then reduce heat and cook uncovered for 15--20 minutes, stirring occasionally, until rich and thickened.
    \item Stir in the frozen vegetables and spinach. Cook for a further 5--10 minutes until the vegetables are tender and the spinach has wilted.
    \item Meanwhile, cook the rigatoni in well-salted boiling water until al dente, then drain.
    \item Off the heat, toss the drained rigatoni through the sauce until well coated.
    \item Scatter mozzarella and Parmesan over the top and stir through while hot until melted, or transfer to an ovenproof dish and grill briefly until bubbling.
  \end{method}
  \begin{notes}
    Originally a deconstructed lasagna concept --- broken lasagna sheets cooked directly in the sauce. This rigatoni version is easier for a weeknight: pasta cooked separately, tossed through at the end.

    If you've already cooked the pasta before making the sauce, keep the extra stock/water in the sauce regardless --- it needs the volume since the pasta won't be absorbing any of it.
  \end{notes}
\end{recipebox}
